\paragraph{Strategy.}
Our scorer reproduces the competition's micro PQ: it gives identical
scores to the organizers' evaluation functions and to an independent dense
re-implementation on random multi-annotator scenes with missed, split,
merged and spurious filaments. We report OOF scores: every
training image is predicted by the models of the fold that did not see its
day. Uncertainty comes from 2{,}000 bootstrap resamples of whole
observations. Calibration is nested: for each outer fold, the 14
reconstruction parameters are chosen by coordinate ascent on the other four
folds (300 OOF images, 332 distinct settings evaluated) and scored on it;
the filter's candidate set, model and threshold are likewise chosen by an
inner cross-validation that never sees the outer fold. The leaderboard was
used only as a sanity check, never for a choice.

\begin{table}[t]
  \caption{OOF results on 707 observations (official micro PQ). Nested:
  every fold scored with settings chosen without it. Tuned: settings chosen
  on all folds (slightly optimistic). T: ConvNeXt-T, S: ConvNeXt-S.}
  \label{tab:main}
  \small
  \setlength{\tabcolsep}{2.5pt}
  \begin{tabular}{lcccc}
    \toprule
    & PQ & SQ & RQ & TP/FP/FN \\
    \midrule
    Human (annotator pairs) & 0.340 & 0.634 & 0.536 & 4292/3709/3709 \\
    T, default post-proc. & 0.416 & 0.675 & 0.616 & 5417/3975/2782 \\
    T + S, default & 0.425 & 0.678 & 0.627 & 5488/3827/2711 \\
    T + S, tuned reconstr. & 0.439 & 0.678 & 0.648 & 5340/2948/2859 \\
    T + S, nested reconstr. & 0.438 & -- & -- & -- \\
    \textbf{T + S, nested + filter} & \textbf{0.447} & -- & -- & -- \\
    \bottomrule
  \end{tabular}
\end{table}

\begin{table}[t]
  \caption{Contribution of each step (OOF PQ on 707 observations, nested
  for the calibrated steps; T: ConvNeXt-T, S: ConvNeXt-S) and public
  leaderboard score of the submissions made with it.}
  \label{tab:steps}
  \small
  \setlength{\tabcolsep}{4pt}
  \begin{tabular}{lcccc}
    \toprule
    & \multicolumn{2}{c}{Nested OOF PQ} & \multicolumn{2}{c}{Leaderboard} \\
    Step & T & T + S & T & T + S \\
    \midrule
    Default post-processing & 0.416 & 0.425 & 0.35 & -- \\
    + calibrated reconstruction & 0.429 & 0.438 & -- & -- \\
    + filter, map features & 0.435 & 0.443 & 0.38 & -- \\
    + filter, map + image feat. & 0.441 & 0.447 & 0.38 & 0.39 \\
    \bottomrule
  \end{tabular}
\end{table}

\paragraph{Results.}
Table~\ref{tab:main} gives the OOF scores and Table~\ref{tab:steps} the
gain of each step. The ensemble improves on ConvNeXt-T alone in all five
folds (default post-processing: 0.422, 0.411, 0.437, 0.421, 0.433 against
0.408, 0.410, 0.432, 0.409, 0.422); we chose it on fold~0 before training
the other folds, where a paired bootstrap over its 139 held-out observations
gave $+0.015$ PQ (95\% CI $+0.007$ to $+0.024$) with calibrated
post-processing, against $+0.010$ ($+0.002$ to $+0.019$) for a second
ConvNeXt-T seed. The calibrated reconstruction mainly removes false
positives (FP $3827\to2948$, RQ $0.627\to0.648$) at constant SQ. The filter
adds $+0.009$ nested PQ and improves every outer fold (0.434$\to$0.441,
0.429$\to$0.439, 0.447$\to$0.460, 0.436$\to$0.443, 0.446$\to$0.452); an
oracle filter on the same candidates would reach 0.511, so ranking
candidates is still the largest source of loss. Fig.~\ref{fig:rubric}
shows the rubric distributions: with the tuned reconstruction, 3.7\% of
ground-truth filaments are fragmented, 3.0\% of predictions merge several
filaments, 16.4\% of ground-truth filaments are missed and 16.2\% of
predictions are spurious; matched pairs have a median IoU of 0.68 (human
pairs: 0.63). Fig.~\ref{fig:test} shows predictions on test images and
Fig.~\ref{fig:curves} the training curves. Predicting the
maps of one test image with 10 models and 4 flips takes 49.5\,s on an
NVIDIA T4; reconstruction and filter take 1.3\,s on 4 CPU cores. The
submission has 1{,}043 filaments on the 180 test images.
Component ablations (clDice, orientation head, endpoint repair, flip TTA)
were planned but not run within the GPU budget; their individual
contribution is therefore not measured.

\begin{figure}[t]
  \centering
  \includegraphics[width=\linewidth]{oof_rubric_panels_default.png}
  \caption{OOF rubric of the T + S ensemble (default post-processing):
  IoU and Dice of overlapping ground-truth/prediction pairs, and
  one-to-one, one-to-many (fragmentation) and many-to-one (merging)
  relations.}
  \label{fig:rubric}
\end{figure}

\begin{figure}[t]
  \centering
  \includegraphics[width=\linewidth]{training_curves.png}
  \caption{Training loss and official PQ of the held-out fold during
  training (default post-processing, no TTA), for the ten models.}
  \label{fig:curves}
\end{figure}

\begin{figure}[t]
  \centering
  \includegraphics[width=\linewidth]{test_predictions.png}
  \caption{Filaments predicted by the submitted pipeline on three test
  images (one color per instance); the elongated filament of the middle
  image is predicted as a single instance.}
  \label{fig:test}
\end{figure}

\paragraph{Strengths and weaknesses.}
The dense centerline and orientation fields keep long, curved filaments in
one piece (fragmentation 3.7\%), and the filter turns the PQ trade-off
between missed and spurious filaments into an explicit, per-instance
decision; both are validated without leakage between days. The main
weakness is the choice of which faint features count as filaments: misses
and spurious predictions are equally frequent (16\%), as between human
annotators (Fig.~\ref{fig:inputs}), and the gap to the oracle filter
(0.447 vs.\ 0.511) shows that instance ranking can still improve. The
public leaderboard (0.39) is below the OOF estimate (0.447), which suggests
that the test images differ from the training days; we did not adapt to it.
Possible improvements: global instead of greedy endpoint matching, more
ensemble members, and self-supervised pre-training on unlabeled GONG frames.
